\clearpage
\SourceLeaf{179}
\pdfbookmark[0]{Göttingische gelehrte Anzeigen. 1831, April 23.}{notice1831april23}
\begin{center}
\rule{.82\linewidth}{.35pt}\par
Göttingische gelehrte Anzeigen.\quad 1831, April 23.\par
\rule{.82\linewidth}{.35pt}
\end{center}
A paper presented to the Royal Society on April 15 by Court Councillor \textsc{Gauss}:
\begin{center}\itshape
Theoria residuorum biquadraticorum, commentatio secunda,
\end{center}
continues the paper already printed in the sixth volume of the \emph{Commentationes novae}, which was also reviewed in these pages at the time, in 1825, on April 11. Even this continuation, although more than twice as long as the first paper, does not yet exhaust this exceptionally rich subject, and completion of the whole will remain reserved for a future third paper.

Although the basic concepts of these theories and the contents of the first paper may be assumed known to everyone who has made higher arithmetic a subject of study, we nevertheless wish to recall those concepts briefly here for the convenience of readers interested in this part of mathematics who do not have the first paper immediately to hand. With respect to any integer $p$, another integer $k$ is called a biquadratic residue if there are numbers of the form $x^4-k$ that are divisible by $p$; in the opposite case it is called a biquadratic nonresidue of $p$. It is sufficient here to restrict attention to the case where $p$ is a prime of the form $4n+1$, and $k$ is not divisible by
\SourceLeaf{180}
that prime, since all other cases are either clear in themselves or can be reduced to this one.

For such a \emph{given} value of $p$, all numbers not divisible by $p$ fall into \emph{four} classes. One contains the biquadratic residues, a second contains those biquadratic nonresidues that are quadratic residues of $p$, and the biquadratic nonresidues that are also quadratic nonresidues are divided between the remaining two. The principle of this division is that either $k^n-1$, or $k^n+1$, or $k^n-f$, or $k^n+f$ will always be divisible by $p$, where $f$ denotes an integer that makes $ff+1$ divisible by $p$. Anyone familiar with the elementary terminology will see how these verbal definitions can be expressed in it.

The theory of this classification is completed in the first paper, both for the immediately evident case $k=-1$ and for the cases $k=\pm2$, which require subtle auxiliary investigations. The beginning of the present paper now proceeds to larger values of $k$. At first, however, only values that are themselves primes need be considered, and the outcome shows that the results take their simplest form when the values are taken positively or negatively according as their absolute values have the form $4m+1$ or $4m+3$. Examining examples immediately yields a rich harvest of new theorems with great ease; we give only a few of them here. The numbering of the classes as $1,2,3,4$ refers to the cases where $k^n$ is congruent to $1,f,-1,-f$, respectively. At the same time, the value always chosen for $f$ is the one that makes $a+bf$ divisible by $p$, where $aa+bb$ represents the decomposition of $p$ into an odd and an even square. Thus examining examples shows that $-3$ always belongs to class $1,2,3,4$ according as $b,a+b,a,a-b$, respectively, is divisible by $3$; that $+5$ belongs to these classes in order according as $b,a-b,a,a+b$, respectively, is divisible by $5$; and that $-7$ falls into class $1$ when $a$ or $b$; into class $2$ when $a-2b$ or $a-3b$; into class $3$ when $a-b$ or $a+b$; and into class $4$ when $a+2b$ or $a+3b$ is divisible by $7$, divisibility being the condition in each of these cases. Similar theorems result for $-11,+13,+17,-19,-23$, and so on. Yet, easy as it is to discover all such particular theorems by examining examples, it seems difficult
\SourceLeaf{181}
to find a general law for these forms by that route, even though some common features soon become apparent; and finding proofs of these theorems is much more difficult still. The methods used for $+2$ and $-2$ in the first paper no longer apply here. Although other methods might also settle what concerns the first and third classes, they prove inadequate for establishing \emph{complete} proofs.

One therefore soon recognizes that this rich region of higher arithmetic can be entered only by entirely new routes. In the first paper, the author had already indicated that a particular extension of the whole field of higher arithmetic was essential for this purpose, without then explaining more closely what it consisted in. The present paper is intended to bring this subject to light.

This consists in nothing other than the following: to give the theory of biquadratic residues its true foundation, the field of higher arithmetic, which had otherwise been extended only to the real integers, must also be extended to the imaginary ones, and these must be granted entirely equal rights of citizenship with the former. Once this is understood, that theory appears in an entirely new light, and its results acquire a most surprising simplicity.

Before the theory of biquadratic residues itself can be developed in this extended domain of numbers, however, the preceding parts of higher arithmetic, which have so far been treated only in relation to real numbers, must share in this extension. We can mention only a few of these preliminary investigations here. The author calls a quantity $a+bi$, where $a$ and $b$ denote real quantities and $i$ is written for $\surd-1$ for brevity, a complex integer when $a$ and $b$ are both integers. Complex quantities thus do not stand opposed to real ones: they include them as the special case where $b=0$. For convenient use, several concepts relating to complex quantities needed particular names, but we shall try to avoid those names in this review.

Just as the arithmetic of real numbers speaks of only two units, the positive and the negative, so in the arithmetic of
\SourceLeaf{182}
complex numbers we have four units, $+1$, $-1$, $+i$, $-i$. A complex integer is called \emph{composite} if it is the product of two integer factors that are not units; a complex number that does not admit \emph{such} a factorization is called a complex prime. For example, the real number $3$ is prime even when considered as a complex number, whereas $5$ is composite as a complex number $=(1+2i)(1-2i)$. Just as in the higher arithmetic of real numbers, primes play a principal role in this extended field of the science.

If a complex integer $a+bi$ is taken as modulus, one can give $aa+bb$ pairwise incongruent complex numbers, and no more, to one of which any given complex integer must be congruent; these can be called a complete system of incongruent residues. The so-called least and absolutely least residues in the arithmetic of real numbers have their exact counterparts here as well. For example, for the modulus $1+2i$, the complete system of absolutely least residues consists of $0$, $1$, $i$, $-1$ and $-i$. Almost all the investigations in the first four sections of the \emph{Disquisitiones Arithmeticae} also have their place, with some modifications, in the extended arithmetic. The celebrated theorem of \textsc{Fermat}, for example, takes the following form here: if $a+bi$ is a complex prime, and $k$ a complex number not divisible by it, then always $k^{aa+bb-1}\equiv1$ for the modulus $a+bi$. Especially remarkable, however, is that the Fundamental Theorem for quadratic residues has its exact counterpart in the arithmetic of complex numbers, in an even simpler form: if $a+bi$, $A+Bi$ are complex primes, with $a$ and $A$ odd, $b$ and $B$ even, then the first is a quadratic residue of the second if the second is a quadratic residue of the first, whereas the first is a quadratic nonresidue of the second if the second is a quadratic nonresidue of the first.

After these preliminary investigations, the paper turns to the theory of biquadratic residues itself. In place of the mere distinction between biquadratic residues and nonresidues, it first establishes a division of the numbers not divisible by the modulus into four classes. If the modulus is a complex prime $a+bi$, with $a$ always assumed odd and $b$ even, and $p$ written for $aa+bb$ for brevity, and if $k$ is a complex number not divisible by $a+bi$, then $k^{\frac14(p-1)}$ will always be
\SourceLeaf{183}
congruent to one of $+1,+i,-1,-i$. This establishes a division of all numbers not divisible by $a+bi$ into four classes, to which the biquadratic characters $0,1,2,3$ are assigned in order. Clearly, character $0$ refers to the biquadratic residues, the others to the biquadratic nonresidues, with character $2$ also corresponding to quadratic\textperiodcentered\ residues, and characters $1$ and $3$ to quadratic nonresidues.

It is easy to see that the main task is to determine this character only for values of $k$ that are themselves complex primes, and here examining examples immediately leads to exceedingly simple results.

If first $k=1+i$, the character of this number is always found to be $\equiv\frac18(-aa+2ab-3bb+1)(\mathrm{mod.}\ 4)$, and similar expressions are found for the cases $k=1-i$, $k=-1+i$, $k=-1-i$.

If, on the other hand, $k=\alpha+\beta i$ is a prime with $\alpha$ odd and $\beta$ even, examining examples very easily yields a law of reciprocity entirely analogous to the Fundamental Theorem for quadratic residues. It can be expressed most simply as follows:

If both $\alpha+\beta-1$ and $a+b-1$ are divisible by $4$ (all other cases can easily be reduced to this one), and the character of $\alpha+\beta i$ with respect to the modulus $a+bi$ is denoted by $\lambda$, whereas the character of $a+bi$ with respect to the modulus $\alpha+\beta i$ is denoted by $l$, then $\lambda=l$ if one of $\beta,b$ (or both) is divisible by $4$. whereas $\lambda=l\pm2$ if neither of $\beta,b$ is divisible by $4$.

These theorems contain essentially everything of substance in the theory of biquadratic residues. Yet, easy as it was to discover them by examining examples, it is difficult to give rigorous proofs of them, especially of the second, the Fundamental Theorem of biquadratic residues. Because of the great length the present paper has already reached, the author was compelled to leave the presentation of the proof of the latter theorem, which he has possessed for $20$ years, for a future third paper. The present paper does, however, give the complete proof of the former theorem concerning $1+i$ (on which the
\SourceLeaf{184}
other theorems for $1-i$, $-1+i$, $-1-i$ depend), and this can already give some idea of the subject's complexity.

We still have a few general remarks to add. Moving the theory of biquadratic residues into the domain of complex numbers might seem objectionable and unnatural to someone less familiar with the nature of imaginary quantities and caught up in false ideas about them. It might suggest that the investigation is thereby left, as it were, hanging in the air, becomes unstable, and moves entirely beyond intuitive understanding. Nothing could be less well founded than such an opinion. On the contrary, the arithmetic of complex numbers admits the clearest visual representation. Although the author has followed a purely arithmetical treatment in this presentation, he has also given the necessary indications for this visual representation, which makes understanding more vivid and is therefore highly to be recommended. these will suffice for readers who think for themselves. Just as absolute integers are represented by a row of equally spaced points on a straight line, in which the starting point represents $0$, the next point $1$, and so on; and just as representing negative numbers then requires only an unbounded extension of this row on the opposite side of the starting point: so representing complex integers requires only the addition that this row be regarded as lying in a fixed unbounded plane, and that an unlimited number of similar rows be taken parallel to it on both sides, at equal distances from one another. In place of a row of points, we then have a system of points that can be arranged as rows of rows in two ways, and that serves to divide the whole plane into equal squares. The point nearest $0$ in the first adjacent row on one side of the row representing the real numbers then corresponds to $i$, just as the point nearest $0$ in the first adjacent row on the other side corresponds to $-i$, and so on. With this representation, the execution of arithmetical operations on complex quantities, congruence, the formation of a complete system of incongruent numbers for a given modulus, and so on, admit a visual representation that leaves nothing to be desired.

\SourceLeaf{185}
On the other hand, this places the true metaphysics of imaginary quantities in a new, bright light.

Our general arithmetic, whose scope far surpasses the geometry of the ancients, is entirely a creation of modern times. Starting originally from the concept of absolute integers, it has extended its domain step by step: fractions were added to integers, irrational numbers to rational numbers, negative numbers to positive numbers, imaginary numbers to real numbers. Yet each advance was initially made with timid, hesitant steps. The first algebraists still called the negative roots of equations false roots, and they are indeed such when the problem to which they relate is presented in a form where the nature of the quantity sought admits no opposite. But just as there is no hesitation about including fractions in \emph{general} arithmetic, although there are so many countable things for which a fraction has no meaning, so negative numbers could not be denied equal rights with positive ones in that arithmetic because countless things admit no opposite: the reality of negative numbers is sufficiently justified, since they find an adequate substrate in countless other cases. All this has, of course, long been understood. But imaginary quantities, set against real ones---formerly. and here and there even now, though inappropriately, called \emph{impossible}---are still more tolerated than naturalized as citizens. They thus appear more like a play of signs empty of content in itself, to which any conceivable substrate is unconditionally denied, without anyone wishing to spurn the rich tribute this play of signs ultimately contributes to the treasury of relations among real quantities.

For many years the author has viewed this highly important part of mathematics from a different standpoint, from which imaginary quantities can be given an object just as well as negative ones. Until now, however, there has been no occasion to state it explicitly in public, although attentive readers will readily find traces of it in the paper on equations published in $1799$ and in the prize essay on the transformation of surfaces. The present paper briefly indicates its main features, which are as follows.

Positive and negative numbers can have an application only where
\SourceLeaf{186}
what is counted has an opposite that, conceived as combined with it, is equivalent to cancellation. Strictly considered, this condition holds only where what is counted consists of relations between pairs of objects, rather than substances (objects conceivable in themselves). It is postulated that these objects are arranged in a row in a definite way, for example $A,B,C,D\dotsfour$, and that the relation of $A$ to $B$ can be regarded as equal to the relation of $B$ to $C$, and so on. Here the concept of opposition requires nothing more than \emph{interchanging} the members of the relation, so that if the relation (or transition) from $A$ to $B$ counts as $+1$, the relation from $B$ to $A$ must be represented by $-1$. Thus, insofar as such a row is unbounded on both sides, every real integer represents the relation of a member chosen arbitrarily as starting point to a definite member of the row.

But if the objects are of a kind that cannot be arranged in a single row, even an unbounded one, and can only be arranged in rows of rows, or, equivalently, form a manifold of two dimensions; and if the relations of one row to another, or the transitions from one to another, behave similarly to the earlier transitions from one member of a row to another member of the same row, then measuring the transition from one member of the system to another clearly requires, besides the former units $+1$ and $-1$, two further units $+i$ and $-i$, likewise opposite to one another. But it must also be postulated that the unit $i$ always denotes the transition from a given member of a row to a \emph{definite} member of the immediately adjacent row. In this way the system can therefore be arranged as rows of rows in two ways.

The mathematician abstracts entirely from the nature of the objects and the content of their relations; the sole concern is to count the relations and compare them with one another. In this respect, just as the relations denoted by $+1$ and $-1$, considered in themselves, are taken to be of the same kind, the mathematician is entitled to extend this to all four elements $+1$, $-1$, $+i$ and $-i$.

These relations can be brought within intuition only by a representation
\SourceLeaf{187}
in space. The simplest case is where there is no reason to arrange the symbols of the objects in any way other than a square pattern: an unbounded plane is divided into squares by two systems of parallel lines crossing at right angles, and the intersections are chosen as the symbols. Each such point $A$ has four neighbours. If the relation of $A$ to one neighbouring point is denoted by $+1$, the relation to be denoted by $-1$ is thereby determined, whereas either of the other two may be chosen for $+i$, or the point corresponding to $+i$ may be taken at will to the \emph{right} or to the \emph{left}. Once forward and backward \emph{in} the plane, and above and below with respect to its two sides, have been fixed (at will), this distinction between right and left is entirely determined \emph{in itself}, although we can communicate our intuition of this distinction to others \emph{only} by pointing to actually existing material things\begingroup\renewcommand{\thefootnote}{*)}\footnote{\textsc{Kant} has already made both observations, but one cannot understand how this perceptive philosopher could believe he found in the first a proof of his opinion that space is \emph{only} the form of our outer intuition, since the second so clearly proves the opposite, and that space must have a real meaning independent of our mode of intuition.}\endgroup. Even after we have decided the latter point, however, we see that it depended on our arbitrary choice which of the two rows crossing at a point we wished to regard as the principal row, and which direction in it as corresponding to positive numbers. We also see that if we wish to take the relation previously treated as $+i$ for $+1$, we must necessarily take the relation previously denoted by $-1$ for $+i$. In mathematical language, this says that $+i$ is a mean proportional between $+1$ and $-1$, or corresponds to the sign $\surd-1$: we deliberately do not say \emph{the} mean proportional, for $-i$ clearly has an equal claim. The possibility of giving $\surd-1$ an intuitive meaning is thus fully justified here, and nothing more is required to admit this quantity into the domain of the objects of arithmetic.

We believed that this brief presentation of the main points of a new theory of so-called imaginary quantities would serve readers interested in mathematics. If this subject has hitherto been viewed from a mistaken standpoint and mysterious obscurity has been found in it,
\SourceLeaf{188}
this is largely attributable to the unsuitable names. Had $+1$, $-1$, $\surd-1$ been called direct, inverse and lateral units, for example, instead of positive, negative and imaginary (or even impossible) units, such obscurity could scarcely have been spoken of. The author has reserved a fuller treatment of the subject, which is really touched on only incidentally in the present paper, for the future. The question why relations among things that form a manifold of more than two dimensions cannot yield further kinds of quantities admissible in general arithmetic will then also receive its answer.
\BookRule{.28\linewidth}

